\chapter{From 2D Classical Ising to 1D Quantum TFIM}
\label{chap:2d-ising-to-1d-tfim}

The previous chapter solved the square-lattice Ising model by rewriting its
row-to-row transfer matrix as a free-fermion operator.  We now sharpen one
particular consequence of that construction: in a strongly anisotropic limit,
the two-dimensional classical Ising transfer matrix becomes the short
imaginary-time evolution operator of the one-dimensional transverse-field Ising
model (TFIM).  This is the cleanest example of the general principle
\begin{equation}
\label{eq:ising-tfim-main-correspondence-schematic}
    d\text{-dimensional quantum system}
    \quad\longleftrightarrow\quad
    (d+1)\text{-dimensional classical statistical system}.
\end{equation}
For the present case, the correspondence is
\begin{equation}
    \text{1D quantum TFIM}
    \quad\longleftrightarrow\quad
    \text{anisotropic 2D classical Ising model}.
\end{equation}
The additional classical direction is Euclidean imaginary time.

The discussion in this chapter has three complementary goals.  First, we derive
precisely how the row-to-row transfer matrix becomes
$\ee^{-\delta\tau H_{\rm TFIM}}$.  Second, we give the explicit dictionary between
classical couplings and quantum Hamiltonian parameters in the conventions of
these notes.  Third, we explain the physical interpretation of the map: spatial
Ising bonds become the quantum Ising interaction, while rare domain walls in the
transfer direction become transverse-field spin flips.

\section{Row transfer matrix and basis convention}
\label{sec:ising-tfim-row-transfer-review}

Consider the anisotropic zero-field classical Ising model on an $L\times M$
square lattice,
\begin{equation}
\label{eq:ising-tfim-classical-partition-function}
    Z_{L,M}(K_x,K_y)
    =
    \sum_{\{s_{j,m}=\pm1\}}
    \exp\left[
        K_x\sum_{j,m}s_{j,m}s_{j+1,m}
        +
        K_y\sum_{j,m}s_{j,m}s_{j,m+1}
    \right],
\end{equation}
with periodic boundary conditions in both directions.  The index $j$ labels the
one-dimensional spatial direction, while $m$ labels the transfer direction.  In
the quantum interpretation, $m$ will become a discretized imaginary-time index.

Following the convention used throughout these notes, a classical row
configuration is identified with a simultaneous eigenstate of the Pauli
operators $\sigma_j^x$:
\begin{equation}
\label{eq:ising-tfim-x-basis-row-state}
    \sigma_j^x\ket{s_1,\ldots,s_L}_x
    =
    s_j\ket{s_1,\ldots,s_L}_x,
    \qquad
    s_j=\pm1 .
\end{equation}
This convention is chosen so that the resulting quantum Hamiltonian is written
in the same form as the TFIM chapter,
\begin{equation}
\label{eq:ising-tfim-quantum-hamiltonian-target}
    H_{\rm TFIM}
    =
    -J\sum_{j=1}^{L}\sigma_j^x\sigma_{j+1}^x
    -h\sum_{j=1}^{L}\sigma_j^z .
\end{equation}
Many books use the rotated convention
$-J\sum_j\sigma_j^z\sigma_{j+1}^z-\Gamma\sum_j\sigma_j^x$.  The two are related
by a global spin rotation.  In these notes the $x$-basis convention is more
natural because the Jordan--Wigner occupation convention uses
$n_j=(1-\sigma_j^z)/2$.

The symmetrized row-to-row transfer matrix is
\begin{equation}
\label{eq:ising-tfim-symmetric-transfer-matrix}
    T
    =
    A_y^L
    \exp\left(
        \frac{K_x}{2}\sum_{j=1}^{L}\sigma_j^x\sigma_{j+1}^x
    \right)
    \exp\left(
        \kappa_y\sum_{j=1}^{L}\sigma_j^z
    \right)
    \exp\left(
        \frac{K_x}{2}\sum_{j=1}^{L}\sigma_j^x\sigma_{j+1}^x
    \right),
\end{equation}
where
\begin{equation}
\label{eq:ising-tfim-kappa-definition}
    \tanh\kappa_y=\ee^{-2K_y},
    \qquad
    A_y=\sqrt{2\sinh(2K_y)} .
\end{equation}
Equivalently,
\begin{equation}
\label{eq:ising-tfim-kappa-log-coth}
    \kappa_y
    =
    \operatorname{arctanh}(\ee^{-2K_y})
    =
    \frac12\log\coth K_y .
\end{equation}
The normalization $A_y^L$ is important for the absolute classical free energy,
but it becomes an additive constant in the quantum Hamiltonian and can be ignored
when discussing normalized correlators and excitation gaps.

\begin{remark}[Why $\sigma^z$ appears in the vertical transfer operator]
In the $\sigma^x$ basis, $\sigma^z$ flips the sign of $s_j$.  Therefore the
one-site matrix
\begin{equation}
    W_{ss'}=\ee^{K_yss'}
\end{equation}
can be represented as $W=A_y\ee^{\kappa_y\sigma^z}$.  A temporal mismatch
$s\neq s'$ is a domain wall in the transfer direction; in the quantum model it is
generated by the transverse-field operator $\sigma^z$.
\end{remark}

\section{Anisotropic Hamiltonian limit}
\label{sec:ising-tfim-anisotropic-limit}

The classical transfer matrix defines an exact operator
\begin{equation}
\label{eq:ising-tfim-exact-effective-hamiltonian}
    H_{\rm eff}
    =
    -\frac{1}{a_\tau}\log T,
\end{equation}
where $a_\tau$ is an arbitrary microscopic time spacing.  For generic finite
$K_x$ and $K_y$, however, $H_{\rm eff}$ is not simply the local nearest-neighbour
TFIM Hamiltonian.  The logarithm of a product of noncommuting many-body
operators generally contains longer-range and higher-body terms.

The local TFIM is obtained in the \emph{strongly anisotropic limit}.  Introduce a
small imaginary-time lattice spacing $\delta\tau$ and set
\begin{equation}
\label{eq:ising-tfim-anisotropic-scaling}
    K_x=\delta\tau J,
    \qquad
    \kappa_y=\delta\tau h,
    \qquad
    \tanh(\delta\tau h)=\ee^{-2K_y} .
\end{equation}
Thus
\begin{equation}
\label{eq:ising-tfim-strong-anisotropy}
    K_x\to0,
    \qquad
    K_y\to\infty,
    \qquad
    K_x\ee^{2K_y}\;\text{fixed}.
\end{equation}
Indeed, for small $\delta\tau$,
\begin{equation}
\label{eq:ising-tfim-ky-asymptotic}
    \ee^{-2K_y}
    =
    \tanh(\delta\tau h)
    =
    \delta\tau h+O(\delta\tau^3),
    \qquad
    K_y
    =
    -\frac12\log(\delta\tau h)+O(\delta\tau^2).
\end{equation}
The horizontal classical coupling becomes weak, while the vertical coupling
becomes strong.  This is exactly what one expects from an imaginary-time path
integral: neighbouring time slices are almost identical, and spin flips in time
are rare events with probability proportional to $\delta\tau h$.

Define
\begin{equation}
\label{eq:ising-tfim-ax-az-definitions}
    X
    =
    J\sum_{j=1}^{L}\sigma_j^x\sigma_{j+1}^x,
    \qquad
    Z
    =
    h\sum_{j=1}^{L}\sigma_j^z .
\end{equation}
Then \cref{eq:ising-tfim-symmetric-transfer-matrix} becomes
\begin{equation}
\label{eq:ising-tfim-symmetric-strang-product}
    T
    =
    A_y^L
    \ee^{\frac{\delta\tau}{2}X}
    \ee^{\delta\tau Z}
    \ee^{\frac{\delta\tau}{2}X} .
\end{equation}
The symmetric Trotter formula gives
\begin{equation}
\label{eq:ising-tfim-strang-expansion}
    \ee^{\frac{\delta\tau}{2}X}
    \ee^{\delta\tau Z}
    \ee^{\frac{\delta\tau}{2}X}
    =
    \exp\left[
        \delta\tau(X+Z)+O(\delta\tau^3)
    \right].
\end{equation}
Since $X+Z=-H_{\rm TFIM}$ in the convention of
\cref{eq:ising-tfim-quantum-hamiltonian-target}, we obtain
\begin{equation}
\label{eq:ising-tfim-transfer-hamiltonian-limit}
    T
    =
    A_y^L
    \exp\left[
        -\delta\tau H_{\rm TFIM}
        +O(\delta\tau^3)
    \right].
\end{equation}
Equivalently,
\begin{equation}
\label{eq:ising-tfim-effective-hamiltonian-expansion}
    -\frac{1}{\delta\tau}\log T
    =
    H_{\rm TFIM}
    -\frac{L}{\delta\tau}\log A_y\,\id
    +O(\delta\tau^2).
\end{equation}
The second term is an additive constant.  It affects the absolute free energy
but not the eigenvectors, normalized expectation values, or excitation gaps.

\begin{proposition}[Hamiltonian limit of the anisotropic Ising transfer matrix]
\label{prop:ising-tfim-hamiltonian-limit}
With the scaling
\begin{equation}
    K_x=\delta\tau J,
    \qquad
    \tanh(\delta\tau h)=\ee^{-2K_y},
\end{equation}
the row-to-row transfer matrix of the anisotropic two-dimensional classical Ising
model satisfies
\begin{equation}
    T
    =
    A_y^L\ee^{-\delta\tau H_{\rm TFIM}+O(\delta\tau^3)},
\end{equation}
where
\begin{equation}
    H_{\rm TFIM}
    =
    -J\sum_j\sigma_j^x\sigma_{j+1}^x
    -h\sum_j\sigma_j^z .
\end{equation}
Therefore the local quantum Hamiltonian is recovered from the logarithm of the
classical transfer matrix up to an additive constant and errors of order
$O(\delta\tau^2)$ in the Hamiltonian density.
\end{proposition}

\section{Partition functions and the continuum-time limit}
\label{sec:ising-tfim-partition-functions}

The classical partition function is
\begin{equation}
\label{eq:ising-tfim-z-tr-tm}
    Z_{L,M}(K_x,K_y)=\Tr T^M .
\end{equation}
Using \cref{eq:ising-tfim-transfer-hamiltonian-limit}, one obtains
\begin{equation}
\label{eq:ising-tfim-z-to-zq-finite-dtau}
    Z_{L,M}(K_x,K_y)
    =
    A_y^{LM}
    \Tr\exp\left[
        -M\delta\tau H_{\rm TFIM}
        +O(M\delta\tau^3)
    \right].
\end{equation}
If
\begin{equation}
\label{eq:ising-tfim-beta-M-dtau}
    \beta=M\delta\tau
\end{equation}
is held fixed while $M\to\infty$ and $\delta\tau\to0$, then
\begin{equation}
\label{eq:ising-tfim-quantum-classical-partition-function}
    Z_{L,M}(K_x,K_y)
    =
    A_y^{LM}
    Z_{\rm TFIM}(\beta)
    \left[1+\text{controlled Trotter corrections}\right],
\end{equation}
where
\begin{equation}
\label{eq:ising-tfim-quantum-partition-function}
    Z_{\rm TFIM}(\beta)=\Tr\ee^{-\beta H_{\rm TFIM}} .
\end{equation}
For local thermodynamic quantities, the logarithm of the correction is of order
$O(\beta L\delta\tau^2)$ for the symmetric decomposition.  Thus the limit is
controlled at fixed system size and fixed inverse temperature, and remains
controlled for thermodynamic densities when the usual local Trotter estimates are
used.

The normalization factor $A_y^{LM}$ gives an additive contribution to the
classical free energy:
\begin{equation}
\label{eq:ising-tfim-free-energy-shift}
    -\log Z_{L,M}
    =
    -LM\log A_y
    -\log Z_{\rm TFIM}(\beta)
    +\cdots .
\end{equation}
This term has no effect on normalized quantum correlation functions.  It is,
however, part of the absolute classical free energy.  One should therefore keep
it when comparing the full classical free-energy density, but drop it when
extracting the quantum Hamiltonian spectrum.

\section{Dictionary of couplings}
\label{sec:ising-tfim-coupling-dictionary}

There are two equivalent ways to present the coupling dictionary.

\subsection{From quantum TFIM to classical Ising}
\label{subsec:ising-tfim-quantum-to-classical}

Given a TFIM Hamiltonian
\begin{equation}
    H_{\rm TFIM}
    =
    -J\sum_j\sigma_j^x\sigma_{j+1}^x
    -h\sum_j\sigma_j^z,
\end{equation}
and a small imaginary-time spacing $\delta\tau$, the corresponding classical
couplings are
\begin{equation}
\label{eq:ising-tfim-quantum-to-classical-dictionary}
    \boxed{
    K_x=\delta\tau J,
    \qquad
    K_y=-\frac12\log\tanh(\delta\tau h)
    }.
\end{equation}
Equivalently,
\begin{equation}
\label{eq:ising-tfim-quantum-to-classical-kappa}
    \kappa_y=\delta\tau h,
    \qquad
    \tanh\kappa_y=\ee^{-2K_y} .
\end{equation}
The classical lattice has $M=\beta/\delta\tau$ rows in the imaginary-time
direction when representing a quantum thermal partition function at inverse
temperature $\beta$.

\subsection{From classical Ising to quantum TFIM}
\label{subsec:ising-tfim-classical-to-quantum}

Given classical couplings $K_x$ and $K_y$, define the dual transfer-direction
coupling
\begin{equation}
\label{eq:ising-tfim-classical-to-quantum-kappa}
    \kappa_y
    =
    \frac12\log\coth K_y .
\end{equation}
If the transfer matrix is close to the identity in the sense that
$K_x\ll1$ and $\kappa_y\ll1$, then it is naturally interpreted as a short-time
quantum evolution operator.  Choosing an imaginary-time spacing $\delta\tau$,
\begin{equation}
\label{eq:ising-tfim-classical-to-quantum-dictionary}
    \boxed{
    J=\frac{K_x}{\delta\tau},
    \qquad
    h=\frac{\kappa_y}{\delta\tau}
    =
    \frac{1}{2\delta\tau}\log\coth K_y
    }.
\end{equation}
Only the dimensionless ratio
\begin{equation}
\label{eq:ising-tfim-ratio-dictionary}
    \boxed{
    \frac{h}{J}
    =
    \frac{\kappa_y}{K_x}
    =
    \frac{1}{2K_x}\log\coth K_y
    }
\end{equation}
is independent of the arbitrary choice of $\delta\tau$.  The overall scale of
$J$ and $h$ is fixed only after one decides how the classical transfer step is to
be measured in imaginary-time units.

\begin{center}
\begin{tabularx}{0.95\linewidth}{@{}L{0.30\linewidth}Y Y@{}}
\toprule
Object & Classical Ising transfer-matrix language & Quantum TFIM language \\
\midrule
Spatial coordinate & Column index $j$ & Spin-chain site $j$ \\
Transfer coordinate & Row index $m$ & Imaginary time $\tau=m\delta\tau$ \\
Row spin $s_{j,m}$ & Eigenvalue of $\sigma_j^x$ & Order-parameter basis variable \\
Horizontal coupling & $K_x=\delta\tau J$ & Ising exchange $J$ \\
Vertical coupling & $\tanh\kappa_y=\ee^{-2K_y}$ & Transverse field $h=\kappa_y/\delta\tau$ \\
Transfer matrix & $T$ & $A_y^L\ee^{-\delta\tau H_{\rm TFIM}}$ \\
Vertical domain wall & $s_{j,m}\neq s_{j,m+1}$ & Spin flip generated by $\sigma_j^z$ \\
\bottomrule
\end{tabularx}
\end{center}

\section{Critical line and self-duality}
\label{sec:ising-tfim-critical-line-self-duality}

The anisotropic square-lattice Ising model is critical on the Onsager curve
\begin{equation}
\label{eq:ising-tfim-onsager-critical-line}
    \sinh(2K_x)\sinh(2K_y)=1 .
\end{equation}
Using
\begin{equation}
\label{eq:ising-tfim-kappa-duality-identity}
    \sinh(2K_y)\sinh(2\kappa_y)=1,
\end{equation}
this condition is equivalent, for positive couplings, to
\begin{equation}
\label{eq:ising-tfim-critical-kx-kappa}
    \kappa_y=K_x .
\end{equation}
Under the Hamiltonian-limit dictionary $K_x=\delta\tau J$ and
$\kappa_y=\delta\tau h$, this becomes
\begin{equation}
\label{eq:ising-tfim-critical-h-equals-j}
    h=J .
\end{equation}
Thus the thermal critical line of the anisotropic two-dimensional classical Ising
model reduces to the zero-temperature quantum critical point of the
one-dimensional TFIM.

This result is also the transfer-matrix form of Kramers--Wannier duality.  The
coupling $\kappa_y$ is the dual of $K_y$, and the critical line is the self-dual
condition
\begin{equation}
\label{eq:ising-tfim-self-dual-condition}
    K_x=\kappa_y .
\end{equation}
The TFIM inherits this as the quantum self-dual condition $J=h$.  On one side of
this point, $J>h$, the quantum ground state is ferromagnetically ordered in the
$x$ direction.  On the other side, $h>J$, the ground state is paramagnetic.  The
classical high-temperature and low-temperature phases become the two quantum
phases separated by a zero-temperature transition.

\begin{remark}[Thermal versus quantum criticality]
The two-dimensional classical transition is driven by thermal fluctuations.  The
one-dimensional TFIM transition at $T=0$ is driven by quantum fluctuations from
the transverse field.  The transfer-matrix correspondence identifies their
Euclidean statistical weights, and therefore their universal critical theory, not
their microscopic dynamical origin.
\end{remark}

\section{Transfer-matrix spectrum and the TFIM dispersion}
\label{sec:ising-tfim-spectrum-dispersion}

The correspondence is sharper than a relation between critical points.  It maps
the full transfer-matrix quasiparticle spectrum to the quantum excitation
spectrum.

In the free-fermion solution of the two-dimensional Ising transfer matrix, the
one-particle transfer dispersion may be written as
\begin{equation}
\label{eq:ising-tfim-onsager-transfer-dispersion}
    \cosh\varepsilon_{\rm tr}(k)
    =
    \cosh(2K_x)\cosh(2\kappa_y)
    -
    \sinh(2K_x)\sinh(2\kappa_y)\cos k .
\end{equation}
The transfer gap closes at $k=0$ when $K_x=\kappa_y$, reproducing the critical
condition above.

Now impose the Hamiltonian-limit scaling
\begin{equation}
    K_x=\delta\tau J,
    \qquad
    \kappa_y=\delta\tau h .
\end{equation}
For small $\delta\tau$,
\begin{align}
\label{eq:ising-tfim-dispersion-expansion}
    \cosh\varepsilon_{\rm tr}(k)
    & =
    1
    +2\delta\tau^2
    \left(J^2+h^2-2Jh\cos k\right)
    +O(\delta\tau^4).
\end{align}
Therefore
\begin{equation}
\label{eq:ising-tfim-transfer-to-quantum-dispersion}
    \varepsilon_{\rm tr}(k)
    =
    \delta\tau\,E_{\rm TFIM}(k)+O(\delta\tau^3),
\end{equation}
where
\begin{equation}
\label{eq:ising-tfim-quantum-dispersion}
    E_{\rm TFIM}(k)
    =
    2\sqrt{h^2+J^2-2Jh\cos k}
    =
    2\sqrt{(h-J\cos k)^2+J^2\sin^2 k}
\end{equation}
is exactly the Bogoliubov quasiparticle dispersion of the TFIM in the convention
of \cref{eq:ising-tfim-quantum-hamiltonian-target}.

Thus transfer-matrix eigenvalue ratios become quantum energy gaps:
\begin{equation}
\label{eq:ising-tfim-eigenvalue-gap-map}
    -\log\frac{\lambda_n}{\lambda_0}
    =
    \delta\tau\,(E_n-E_0)+O(\delta\tau^3).
\end{equation}
At the ferromagnetic critical point $h=J>0$,
\begin{equation}
\label{eq:ising-tfim-critical-linear-dispersion}
    E_{\rm TFIM}(k)
    =
    4J\abs{\sin(k/2)}
    \simeq
    2J\abs{k}
    \qquad
    (k\to0).
\end{equation}
The continuum limit is therefore a massless Majorana theory.  Away from the
critical point, the quantum gap is
\begin{equation}
\label{eq:ising-tfim-quantum-gap}
    \Delta
    =
    E_{\rm TFIM}(0)
    =
    2\abs{h-J}
\end{equation}
for $J,h>0$.

\section{Operator dictionary and Euclidean correlators}
\label{sec:ising-tfim-operator-dictionary}

The transfer-matrix map also identifies correlation functions.  A classical spin
on row $m$ becomes the quantum operator $\sigma_j^x$ inserted at imaginary time
\begin{equation}
    \tau=m\delta\tau .
\end{equation}
Thus, in the Hamiltonian limit,
\begin{equation}
\label{eq:ising-tfim-spin-correlator-map}
    \langle s_{j,m}s_{j',0}\rangle_{\rm cl}
    \longrightarrow
    \frac{1}{Z_{\rm TFIM}}
    \Tr\left[
        \ee^{-\beta H_{\rm TFIM}}
        \sigma_j^x(\tau)\sigma_{j'}^x(0)
    \right],
\end{equation}
where
\begin{equation}
\label{eq:ising-tfim-imaginary-time-heisenberg}
    \sigma_j^x(\tau)
    =
    \ee^{\tau H_{\rm TFIM}}\sigma_j^x\ee^{-\tau H_{\rm TFIM}} .
\end{equation}
At zero quantum temperature, $\beta\to\infty$, the trace becomes a ground-state
expectation value:
\begin{equation}
\label{eq:ising-tfim-zero-temperature-correlator}
    \langle s_{j,m}s_{j',0}\rangle_{\rm cl}
    \longrightarrow
    \bra{0}\sigma_j^x(\tau)\sigma_{j'}^x(0)\ket{0} .
\end{equation}

Spatial classical bonds map to the Ising energy operator,
\begin{equation}
\label{eq:ising-tfim-spatial-bond-map}
    s_{j,m}s_{j+1,m}
    \quad\longleftrightarrow\quad
    \sigma_j^x\sigma_{j+1}^x .
\end{equation}
Temporal bonds are represented by insertions obtained by differentiating the
vertical transfer matrix with respect to $K_y$.  Since
\begin{equation}
    W=A_y\ee^{\kappa_y\sigma^z},
\end{equation}
vertical-bond observables contain an identity contribution and a contribution
proportional to $\sigma^z$.  In the quantum Hamiltonian limit, the nontrivial
operator associated with temporal fluctuations is therefore the transverse-field
operator
\begin{equation}
\label{eq:ising-tfim-temporal-bond-map}
    \text{temporal domain-wall fluctuations}
    \quad\longleftrightarrow\quad
    \sigma_j^z .
\end{equation}
This is the operator-level reason why the vertical classical coupling becomes
the transverse field rather than another Ising exchange term.

\section{Physical interpretation of the anisotropic limit}
\label{sec:ising-tfim-physical-interpretation}

The anisotropic limit has a simple worldline interpretation.  In the $\sigma^x$
basis, a quantum spin configuration at imaginary time $\tau$ is a row
configuration
\begin{equation}
    \bm{s}(\tau)=(s_1(\tau),\ldots,s_L(\tau)).
\end{equation}
The diagonal part of the Hamiltonian,
\begin{equation}
    -J\sum_j\sigma_j^x\sigma_{j+1}^x,
\end{equation}
assigns an energy to each row configuration.  Over a short imaginary-time
interval $\delta\tau$, it contributes the Boltzmann factor
\begin{equation}
\label{eq:ising-tfim-diagonal-row-weight}
    \exp\left[
        \delta\tau J\sum_j s_j(\tau)s_{j+1}(\tau)
    \right],
\end{equation}
which is precisely the horizontal classical Ising weight with
$K_x=\delta\tau J$.

The off-diagonal part of the Hamiltonian,
\begin{equation}
    -h\sum_j\sigma_j^z,
\end{equation}
flips $\sigma^x$ eigenvalues.  For one spin,
\begin{equation}
\label{eq:ising-tfim-one-site-short-time-matrix}
    \bra{s'}_x\ee^{\delta\tau h\sigma^z}\ket{s}_x
    =
    \begin{cases}
    \cosh(\delta\tau h), & s'=s,\\
    \sinh(\delta\tau h), & s'=-s.
    \end{cases}
\end{equation}
The relative weight of a flip is therefore
\begin{equation}
\label{eq:ising-tfim-flip-weight}
    \frac{\sinh(\delta\tau h)}{\cosh(\delta\tau h)}
    =
    \tanh(\delta\tau h)
    =
    \ee^{-2K_y}.
\end{equation}
A vertical classical domain wall costs a factor $\ee^{-2K_y}$ relative to a
non-domain-wall segment.  Hence the transverse field is the rate per unit
imaginary time at which a worldline flips:
\begin{equation}
\label{eq:ising-tfim-transverse-field-as-flip-rate}
    \ee^{-2K_y}\simeq\delta\tau h .
\end{equation}
This is the microscopic origin of the dictionary.

The ordered and disordered phases can be read from this worldline picture:
\begin{itemize}[leftmargin=2.0em]
    \item If $J\gg h$, spatial alignment dominates and spin worldlines remain
    ordered in the $x$ direction.  This is the ferromagnetic phase of the TFIM.
    \item If $h\gg J$, spin flips proliferate in imaginary time and destroy
    long-range $x$ order.  This is the quantum paramagnet.
    \item At $h=J$, domain-wall proliferation is critical.  The transfer-matrix
    gap closes and the Euclidean theory becomes scale invariant.
\end{itemize}

\section{Finite temperature, finite size, and the extra dimension}
\label{sec:ising-tfim-finite-temperature-finite-size}

The number of rows $M$ in the transfer direction determines the quantum inverse
temperature:
\begin{equation}
    \beta=M\delta\tau .
\end{equation}
Therefore:
\begin{itemize}[leftmargin=2.0em]
    \item A classical cylinder with finite circumference $M$ in the transfer
    direction represents a quantum TFIM at finite temperature $T_{\rm q}=1/\beta$.
    \item The limit $M\to\infty$ at fixed $\delta\tau$ projects onto the largest
    eigenvalue of the transfer matrix.  In the quantum language this is the
    zero-temperature or ground-state limit.
    \item The continuum-time quantum limit requires $M\to\infty$ and
    $\delta\tau\to0$ with $\beta=M\delta\tau$ fixed.
\end{itemize}
The spatial size $L$ is the number of sites in the quantum spin chain.  Thus the
classical thermodynamic limit $L,M\to\infty$ contains both the quantum
thermodynamic limit and the zero-temperature limit.

The transfer-matrix correspondence also explains why a quantum critical point in
$d$ spatial dimensions is naturally described by a $(d+1)$-dimensional Euclidean
field theory.  The quantum partition function
\begin{equation}
    Z_{\rm q}=\Tr\ee^{-\beta H}
\end{equation}
is a sum over histories in imaginary time.  When the quantum gap closes, the
correlation length in space and the correlation length in imaginary time both
diverge.  The relative scaling of the two is controlled by the dynamical exponent
$z$.  For the TFIM critical point, $z=1$, and the long-distance continuum theory
is effectively isotropic after a nonuniversal rescaling of time.

\section{Common pitfalls}
\label{sec:ising-tfim-common-pitfalls}

There are several points that are easy to misstate.

\begin{enumerate}[label=(\roman*)]
    \item \textbf{The generic transfer matrix is not automatically a local
    TFIM Hamiltonian.}  For finite $K_x$ and $K_y$, the exact logarithm
    $-\log T$ is a well-defined free-fermion Hamiltonian, but it is generally not
    the nearest-neighbour TFIM.  Locality in the simple TFIM form emerges in the
    strongly anisotropic limit.

    \item \textbf{The normalization $A_y^L$ is not a physical transverse field.}
    It shifts the absolute free energy or the zero of quantum energy.  The
    transverse field is encoded in $\kappa_y$, not in $A_y$.

    \item \textbf{The choice of spin basis matters.}  In these notes the
    classical Ising variable is the eigenvalue of $\sigma^x$.  If one instead
    labels classical spins by $\sigma^z$, the same physics is obtained after a
    global spin rotation, but the symbols in the Hamiltonian are interchanged.

    \item \textbf{Classical temperature and quantum temperature are different
    concepts.}  The classical couplings $K_x,K_y$ already include the classical
    inverse temperature.  The quantum inverse temperature is the length of the
    transfer direction, $\beta=M\delta\tau$.

    \item \textbf{Boundary conditions introduce fermion parity sectors.}  After
    Jordan--Wigner transformation, periodic spin boundary conditions split the
    free-fermion solution into parity-dependent boundary conditions.  This
    sector bookkeeping is essential for exact finite-size spectra, although it
    does not change the local Hamiltonian-limit dictionary above.
\end{enumerate}

\section{Summary}
\label{sec:ising-tfim-summary}

The anisotropic two-dimensional classical Ising model contains the
one-dimensional quantum TFIM as its Hamiltonian limit.  The essential chain of
relations is
\begin{equation}
\label{eq:ising-tfim-summary-chain}
\begin{gathered}
    Z_{\rm 2D\ Ising}=\Tr T^M,
    \qquad
    T
    =
    A_y^L
    \ee^{\frac{K_x}{2}\sum_j\sigma_j^x\sigma_{j+1}^x}
    \ee^{\kappa_y\sum_j\sigma_j^z}
    \ee^{\frac{K_x}{2}\sum_j\sigma_j^x\sigma_{j+1}^x},
    \\
    K_x=\delta\tau J,
    \qquad
    \kappa_y=\delta\tau h,
    \qquad
    \tanh\kappa_y=\ee^{-2K_y},
    \\
    T
    =
    A_y^L\ee^{-\delta\tau H_{\rm TFIM}+O(\delta\tau^3)},
    \qquad
    H_{\rm TFIM}
    =
    -J\sum_j\sigma_j^x\sigma_{j+1}^x
    -h\sum_j\sigma_j^z .
\end{gathered}
\end{equation}
The coupling ratio is
\begin{equation}
\label{eq:ising-tfim-summary-ratio}
    \frac{h}{J}
    =
    \frac{\kappa_y}{K_x}
    =
    \frac{1}{2K_x}\log\coth K_y .
\end{equation}
The Onsager critical curve becomes the TFIM quantum critical point:
\begin{equation}
\label{eq:ising-tfim-summary-critical}
    \sinh(2K_x)\sinh(2K_y)=1
    \quad\Longleftrightarrow\quad
    K_x=\kappa_y
    \quad\Longleftrightarrow\quad
    J=h .
\end{equation}
At the spectral level,
\begin{equation}
\label{eq:ising-tfim-summary-dispersion}
    \varepsilon_{\rm tr}(k)
    =
    \delta\tau\,2\sqrt{h^2+J^2-2Jh\cos k}
    +O(\delta\tau^3),
\end{equation}
so the transfer-matrix gap is the quantum excitation gap measured in imaginary-
time units.  This chapter therefore gives the canonical solvable example of how
classical transfer matrices, quantum Hamiltonians, duality, and free-fermion
criticality are different representations of the same underlying structure.
